\emph{প্রমাণপদ্ধতি} দিয়ে বাক্যগুলির মধ্যে অনুসিদ্ধান্ত ও
সামঞ্জস্যকে বিশুদ্ধভাবে সংকেতবিন্যাসের ভিত্তিতে চরিত্রায়িত করা
যায়। \emph{সিকোয়েন্ট কলন} এমনই একটি প্রমাণপদ্ধতি। এতে একটি
\emph{নিষ্পাদন} সিকোয়েন্টগুলির একটি বৃক্ষ নিয়ে গঠিত (একটি
সিকোয়েন্ট $\Gamma \Sequent \Delta$ হলো~$\Sequent$ দিয়ে পৃথক
করা সূত্রের দুটি অনুক্রম)। !!a{derivation}-এর একেবারে উপরের
সিকোয়েন্টগুলি $!A \Sequent !A$ আকারের \emph{প্রারম্ভিক সিকোয়েন্ট}।
!!{derivation} সঠিক হতে হলে বাকি প্রত্যেকটি সিকোয়েন্টকে কয়েকটি
\emph{অনুমানবিধির} একটির সাহায্যে যথাযথভাবে প্রতিষ্ঠা করতে হয়।
এই বিধিগুলি জোড়ায় জোড়ায় আসে। প্রতিটি সংযোজক ও পরিমাণসূচকের
জন্য সিকোয়েন্টের বাঁ পাশে প্রয়োগের একটি বিধি ও ডান পাশে
প্রয়োগের একটি বিধি থাকে। যেমন, $\Gamma \Sequent \Delta, !A \lif !B$
সিকোয়েন্টটি $\RightR{\lif}$ বিধি দিয়ে প্রতিষ্ঠা করা হলে
আগের সিকোয়েন্টটি (\emph{পূর্বধারণা}) অবশ্যই
$!A, \Gamma \Sequent \Delta, !B$ হবে। কিছু বিধি সিকোয়েন্টে
বাক্যের ক্রম বা সংখ্যা বদলাতেও দেয়। যেমন, \RightR{\Exchange}
বিধিতে সিকোয়েন্টের ডান পাশের দুটি সূত্রের স্থান অদলবদল করা যায়।

$\quad \Sequent !A$ সিকোয়েন্টটির !!a{derivation} থাকলে
বলি $!A$ একটি \emph{উপপাদ্য} এবং লিখি~$\Proves !A$।
$\Gamma_0 \Sequent !A$-এর !!a{derivation} থাকলে, যেখানে
$\Gamma_0$-এর প্রত্যেক $!B$~$\Gamma$-র অন্তর্ভুক্ত, বলি $!A$ হলো
~$\Gamma$ \emph{থেকে নিষ্পাদনযোগ্য} এবং লিখি $\Gamma \Proves !A$।
$\Gamma_0 \Sequent \quad$-এর !!a{derivation} থাকলে, যেখানে
$\Gamma_0$-এর প্রত্যেক $!B$~$\Gamma$-র অন্তর্ভুক্ত, বলি $\Gamma$
\emph{অসঙ্গত}, অন্যথায় \emph{সঙ্গত}। এই ধারণাগুলি পরস্পর
সম্পর্কিত। যেমন, $\Gamma \Proves !A$ হয় ঠিক তখনই, যখন
$\Gamma \cup \{\lnot !A\}$ অসঙ্গত। অনুরূপ অর্থগত ধারণাগুলির সঙ্গেও
এদের সম্পর্ক আছে। যেমন, $\Gamma \Proves !A$ হলে $\Gamma \Entails !A$।
প্রমাণপদ্ধতির এই ধর্মকে \emph{বিশুদ্ধতা} বলে—$\Gamma$ থেকে যা
নিষ্পাদন করা যায়, তা~$\Gamma$ থেকে অবশ্যই অনুসৃত হয়।
!!{derivation}s-এর দৈর্ঘ্যের উপর আরোহ করে \emph{বিশুদ্ধতা
উপপাদ্য} প্রমাণ করা হয়। তাতে দেখানো হয় যে প্রতিটি অনুমানে
পূর্বধারণা-সিকোয়েন্টগুলি বৈধ হলে সিদ্ধান্ত-সিকোয়েন্টটিও বৈধ হয়।
